
Введём обозначение $p_j$ - $j$-е простое число. Пусть $dp_{n,j}$ — это количество $k$ таких, что $1 \leq k \leq n$ и все простые делители $k$ не меньше $p_j$ (заметим, что число $1$ будет учтено во всех $dp_{n,j}$, поскольку оно не имеет простых делителей). $dp_{n, j}$ удовлетворяет реккурентному соотношению:

$dp_{n,1} = n$ (поскольку $p_1=2$).

$dp_{n, j} = dp_{n, j} + 1 + dp_{\left[ \frac{n}{p_j}\right], j}$, следовательно $dp_{n, j + 1} = dp_{n, j} - dp_{\left[\frac{n}{p_j}\right], j}$.

Пусть $p_k$ — это наименьшее простое большее $\sqrt{n}$. Тогда $\pi(n) = dp_{n, k} + k - 1$ (по определению первое слагаемое учитывает все простые не меньшие $k$).

Если вычислять реккурентность $dp_{n, k}$ напрямую, то все достижимые состояния будут иметь вид $dp_{\left[ \frac{n}{i}\right], j}$. Также можно заметить, что если $p_j$ и $p_k$ оба больше $\sqrt{n}$, то $dp_{n, j} + j = dp_{n, k} + k$. Поэтому, для всех $\left[\frac{n}{i}\right]$ нам достаточно хранить только $\approx\pi(\sqrt{n/i})$ значений $dp_{\left[ \frac{n}{i}\right], j}$.

Вместо прямого подсчёта всех состояний динамики, будем осуществлять двухшаговый процесс:

Выберем $K$.

Запустим рекурсивное вычисление $dp_{n, k}$. Если при этом в какой-то момент мы захотим посчитать значение для состояние $n < K$, запомним запрос ``посчитать количество чисел не больше $n$ с простыми делителями не меньше $k$''.

Посчитаем ответы на все запросы в off-line: вычислим решето для чисел до $K$, отсортируем все числа по наименьшему простому делителю. Теперь мы можем ответить на все запросы, используя структуру данных на прибавление в точке и взятие суммы на отрезке (например, дерево Фенвика). Запомним все ответы глобально.

Снова запустим рекурсивное вычисление $dp_{n, j}$. Но в этот раз если мы захотим посчитать значение для состояния $n < K$, мы можем использовать уже вычисленное значение.

Эффективность этой идеи сильно зависит от величины $Q$ — количество запросов, которые нам придётся обработать.


Предпосчёт ответов для $Q$ запросов может быть выполнен за время $(K + Q)\log(K + Q)$ и это наиболее ёмкая часть вычислений. $K\approx\left(\frac{n}{\log n}\right)^{2/3}$.
